% Introduction and Related Work draft — literature-agent output.
% Writing agent should rewrite and adapt this in a unified voice.

\section{Introduction}

Bone is a dynamically remodeled organ whose mass depends on a balance between osteoblastic formation and osteoclastic resorption, and accumulating genetic and pharmacological evidence now situates the sympathetic nervous system (SNS) as an active regulator of that balance. Stimulation of the \(\beta_2\)-adrenergic receptor on osteoblasts suppresses bone formation while promoting resorption, and genetic ablation of adrenergic signaling raises bone mass in the mouse skeleton \cite{takeda2002leptin,elefteriou2005leptin}. Leptin, acting centrally through hypothalamic relays, produces an antiosteogenic signal that is propagated to the periphery via the SNS \cite{ducy2000leptin,flak2016cns}. Peripheral consequences of SNS drive at the bone surface have also been characterized: noradrenergic fibers immunoreactive for tyrosine hydroxylase, dopamine \(\beta\)-hydroxylase, or neuropeptide Y densely innervate periosteum and marrow, with vasoactive intestinal peptide-containing fibers representing a distinct parenchymal population \cite{hohmann1986innervation,togari2010control}. Complementary loss-of-function work, such as the osteopenia seen in dopamine-transporter-deficient mice, underscores that catecholaminergic tone itself, not only its receptor signaling, is skeletal-relevant \cite{bliziotes2000bone,zhu2018cortical}.

Despite this physiological evidence, the central architecture that delivers SNS drive to bone in the mouse remains largely uncharted. The alpha-herpesvirus pseudorabies virus (PRV), and in particular the attenuated Bartha-derived reporter strain PRV152, has become the standard retrograde transneuronal tracer for this question: it is taken up at peripheral axon terminals, replicates in post-synaptic somata, and then crosses synapses strictly in the retrograde direction, chaining along functionally connected circuits \cite{card2001transneuronal,ekstrand2008alpha,card1995influence}. Using PRV152 and related reagents, several groups have mapped central SNS outflow to white and brown adipose tissue, defining shared and depot-specific brain nodes \cite{bamshad1998central,bamshad1999cns,bartness1998innervation,bartness2014neural,ryu2014short,song2008mc4r,shrestha2010central,brito2007differential,adler2012neurochemical}. A hierarchical rat bone-marrow circuit that transits through lower thoracic and upper lumbar intermediolateral column (IML) preganglionic neurons and paravertebral chain ganglia has also been described for the rat femoral epiphysis, and a direct neuroanatomical link between phosphodiesterase-5A (PDE5A)-containing neurons in defined brain sites and bone was recently established using PRV152 \cite{kim2020pde5a}.

A map of comparable resolution for the mouse femur has, however, been missing. This leaves two important questions open. First, does the central SNS control of bone simply overlap with that of adjacent peripheral targets (adipose depots, bone marrow), or does the femur recruit a partially private network? Second, do descending pain-modulatory structures such as the periaqueductal gray (PAG) and the medullary raphe contribute to femur-directed SNS outflow, which would give a neuroanatomical substrate for central modulation of bone pain \cite{francois2017brainstem,heinricher2009descending}? Here we close this gap. By injecting PRV152 into the metaphysis and periosteum of the right femur of adult male mice and performing serial immunofluorescent analysis of the neuroaxis, we build an atlas of central SNS outflow to bone. The resulting data span 87 PRV152-positive brain nuclei, sub-nuclei, and regions across six brain divisions, identify hypothalamic, midbrain/pontine, and medullary hubs, and nominate a raphe magnus-periaqueductal gray axis as a candidate relay for bone pain.

Our contributions are: (i) the first comprehensive catalog of PRV152-labeled SNS neurons projecting to the mouse femur; (ii) per-division counts with SEMs that quantify the relative contribution of each brain division to bone-directed SNS outflow; (iii) identification of dominant nuclei within each division, including LH, PVH and DM in the hypothalamus, PAG, LPAG and PnO in the midbrain and pons, and RPa, RMg and Gi in the medulla; (iv) anatomical evidence consistent with a shared-with-site-specific-variation model of SNS control across bone, marrow and adipose tissues; and (v) a candidate circuit diagram for central modulation of bone pain. Together these data provide a neuroanatomical foundation for dissecting autonomic regulation of bone metabolism and for targeting descending pathways to manage bone pain.

\section{Related Work}

\textbf{Peripheral SNS control of bone.} That adrenergic signaling tunes bone remodeling has been repeatedly demonstrated in mouse genetic and pharmacological work \cite{ducy2000leptin,takeda2002leptin,elefteriou2005leptin}. Leptin circulates in proportion to adiposity and acts on hypothalamic LepRb-expressing neurons to engage sympathoexcitation, driving both reduced formation through \(\beta_2\)-adrenergic receptors on osteoblasts and enhanced resorption through RANKL \cite{elefteriou2005leptin,flak2016cns}. Convergent evidence from dopamine-transporter-null mice, which are osteopenic, and from norepinephrine transporter studies in cortical bone, broadens the catecholaminergic picture \cite{bliziotes2000bone,zhu2018cortical}. The peripheral ``endpoint'' of these signals is a dense, heterogeneous innervation of periosteum and marrow that mixes noradrenergic and peptidergic fibers \cite{hohmann1986innervation,togari2010control}.

\textbf{Transneuronal tracing of autonomic circuits.} The attenuated Bartha-derived strains of PRV, including PRV152, have been the workhorses for mapping multi-synaptic SNS circuits from peripheral tissues to the brain \cite{card2001transneuronal,ekstrand2008alpha,card1995influence}. Applied to adipose depots, PRV152 has defined the central SNS outflow to white and brown adipose tissue, with shared nodes in the medial preoptic area, paraventricular hypothalamus, and lateral hypothalamus, and depot-specific biases across suprachiasmatic, dorsomedial, arcuate, and retrochiasmatic regions \cite{bamshad1998central,bamshad1999cns,bartness1998innervation,bartness2014neural,song2008mc4r,shrestha2010central,brito2007differential}. Dual anterograde/retrograde tracing has further resolved short spinal and long brain-level sympathetic-sensory feedback loops relevant to lipolysis \cite{ryu2014short}, and sexually dimorphic projection patterns to rat adipose depots \cite{adler2012neurochemical}. These studies establish the anatomical primitives we leverage for bone.

\textbf{Bone marrow and bone PRV work.} Prior PRV mapping of rat femoral bone marrow identified preganglionic IML neurons at T4-L1, paravertebral chain ganglia, and a limited hierarchical brain network. A direct link between PDE5A-containing brain neurons (in the locus coeruleus, raphe pallidus, and paraventricular hypothalamic nucleus) and bone was recently established using PRV152 together with pharmacology and behavior \cite{kim2020pde5a}. What has been absent is a comprehensive, mouse femur-specific atlas that would enable quantitative comparisons across divisions.

\textbf{Descending pain modulation.} The periaqueductal gray and the rostral ventromedial medulla, including raphe magnus and raphe pallidus, form the canonical descending pain-modulatory system that gates nociceptive input at the dorsal horn \cite{heinricher2009descending,francois2017brainstem,benarroch2012postural}. These same structures also contribute to autonomic control. Our data place bone-directed SNS neurons within these same nuclei, making the anatomical case that bone pain and bone metabolism share partly overlapping descending substrates.
